% !TEX root = ../print.tex
% Section 5.8 of the article: shape-adapted smoothing across the rungs (notes/PLAN-article-structure.md,
% "The outline"). The section is prose of the article, reading the practice of Lindeberg and Garding
% against the families of 5.4-5.6; the prose closing this part, the draft's Remark 2.17 read in the
% nodes' terms, is its material (it closed 05-smaller-groups.tex until Q-0284, the author's decision
% of 2026-10-04). The one node here is the linear algebra its first paragraph needs, as a corollary of
% thm:rank-deficient-step (Q-0282 item 4, the draft rewrite's gap in its 5.8, Q-0279). It is placed in
% its own part, after 5.7, so that the section numbers follow the article's, and so that it does not
% touch 04-classical.tex. Spends nothing.

\section{Shape-adapted smoothing across the rungs}
\label{sec:shape-adapted}

The practice of shape-adapted smoothing normalizes the adapted scale matrix by its least eigenvalue.
The step from the isotropic first stage to the adapted aperture is then rank-deficient, and only the
Gaussian has a positive increment along it (\ref{prop:shape-adapted-apertures}).

% not in the draft as a node: the draft rewrite's 5.8, first paragraph (Q-0279), which applies
% thm:rank-deficient-step to the apertures Sigma_t = t M_I of Lindeberg-Garding 1994, step 2
% (records/affine/READING-lindeberg1994shape.md, section 3: M_I = M / lambda_min(M), local scale
% matrix t M_I). The node states the linear algebra and the corollary on kernels; the source is named
% in the annotation as the reason for the normalization, and no proof reads it.
% CHANGED Q-0281: typed (the statement skeleton); no split (records/formalization/PRICING.md).
% CHANGED Q-0327: notation only (whole-paper review flag 21; blind draft review flag 29): the variance
% of the aperture is $w$, no longer $t$, which is the volume scale of lem:loewner-scale-shape.
% CHANGED round-2 review (Gemini), item 15 (records/affine/reviews/VERIFY-round2-notation-paper.md):
% wording only, "Its kernel" -> "Its null space ker P" in (1), as the paper writes it elsewhere, since
% "kernel" names a law in this paper. No claim changes; \uses and \lean unchanged.
\begin{proposition}[the apertures of shape-adapted smoothing are rank-deficient steps]
\label{prop:shape-adapted-apertures}
\uses{thm:rank-deficient-step,prop:affine-orbits,def:gw-cascade-family,lem:gw-kernels-from-origin,prop:gw-ray-families}
\lean{AffineHemigroup.shape_adapted_step, AffineHemigroup.shape_adapted_gaussian,
AffineHemigroup.shape_adapted_scheme}\leanok
Let $d \ge 2$, $w > 0$, and $M \in \PD(d)$ with least eigenvalue $1$, and put $\Sigma_w := wM$, so that
the least eigenvalue of $\Sigma_w$ is $w$, and $P := \Sigma_w - wI$.
\begin{enumerate}
\item \emph{(The step.)} $P = w(M - I)$ is positive semidefinite. Its null space $\ker P$ is the eigenspace of $M$ for
the eigenvalue $1$, and its range is spanned by the eigenvectors of $M$ with eigenvalues above $1$. So
$\rank P \le d - 1$, with $\rank P = 0$ exactly when $M = I$; for $d = 2$ and $M \ne I$, $\rank P = 1$.
\item \emph{(Only the Gaussian follows the apertures.)} Let $g$, $\mu_I$ and the affine orbit
$\mu_\Sigma = (\Sigma^{1/2})_*\mu_I$, $\Sigma \in \PD(d)$, be as in \ref{prop:affine-orbits}, and let
$M \ne I$. There is a probability measure $\rho$ on $\RR^d$ with $\mu_{\Sigma_w} = \mu_{wI} * \rho$ if and
only if $g(u) = bu$ for some $b \ge 0$, that is, $\mu_I$ is the centred Gaussian law with covariance $2bI$
($\delta_0$ when $b = 0$). Then $\rho$ is the centred Gaussian law with covariance $2bP$, carried by the
range of $P$.
\item \emph{(In the scheme.)} Let $(\GL(d), W)$ be a wedge datum (\ref{def:gw-cascade-family}) with
$wI \le_W wM$ for some $w > 0$ and some $M \ne I$ as above. Then for every $(\GL(d), W)$-cascade family
of kernels there is a $b > 0$ with $\widehat\mu_{\Sigma,\Sigma'}(\omega) =
e^{-b\langle\omega,(\Sigma'-\Sigma)\omega\rangle}$ for every pair $\Sigma \le_W \Sigma'$.
\end{enumerate}
\statusT\quad\emph{(Cites nothing beyond the spectral theorem and the uniqueness of the Fourier
transform. \textbf{What it is worth.} It is the statement the article's reading of shape-adapted
smoothing rests on. The iteration of Lindeberg and Gårding normalizes the second-moment matrix by its
least eigenvalue and uses $wM_I$ as the local scale matrix, $M_I := M/\lambda_{\min}(M)$ (Lindeberg and
Gårding 1994, step 2 of the composed method, p.~396, read at the page:
\texttt{records/affine/READING-lindeberg1994shape.md}); so its apertures are the $\Sigma_w$ of this node,
with least eigenvalue the scale $w$ of the isotropic first stage. By (1) the step $wI \to \Sigma_w$
leaves the least eigendirection unchanged and widens the others, and by (2) such a step is a further
smoothing only in the Gaussian orbit, by the degenerate Gaussian along the widened directions, in the
plane along the major axis alone. In every other affine orbit the adapted aperture is related to the first
stage by no increment that is a probability measure, though the orbit may cascade along other steps of
full rank (\ref{prop:stable-step-cone}). Clause (3) reads it in the scheme: an order invariant under the
general linear group that contains one such step leaves only Gaussian increments, as the Loewner order
does (\ref{thm:affine-rigidity}); the stable orders $\preceq_\alpha$ contain none
(\ref{prop:stable-step-cone}(3)). \textbf{Not claimed}: anything about the convergence of the iteration or
the second-moment matrix; the selection at fixed volume, which no order relates
(\ref{lem:loewner-scale-shape}(2)); and the existence of a $(\GL(d), W)$-family in (3), which needs
$\Sigma' - \Sigma$ positive semidefinite on every pair of $W$.)}
\end{proposition}

% CHANGED Q-0281 wave 1: proof of (1) rewritten to the Lean route: no eigenbasis of M is chosen; the
% least eigenvalue is attained at a unit eigenvector, the range is identified through three facts and the
% spectral theorem in the form "the eigenspaces span R^d", and the rank through rank-nullity.
% CHANGED Q-0328: four semicolons and a trailing "being" clause made sentences (register; no step
% changed), and the paper's copy re-transcribed verbatim (it still printed the eigenbasis argument).
% CHANGED round-2 review (Gemini), items 8 and 15 (records/affine/reviews/VERIFY-round2-structure-paper.md
% § 8, VERIFY-round2-notation-paper.md § 15): the first sentence, which said that the least eigenvalue
% is the least eigenvalue, now reads the hypothesis through the Rayleigh quotient; "the kernel of P" and
% "its kernel" read "the null space ker P" and "its null space". The route (the Lean's, Q-0281) and
% every step are unchanged.
\begin{proof}
\leanok
\emph{(1).} The least eigenvalue of $M$ is $1 = \min_{|x| = 1}\langle x, Mx\rangle$, and it is attained at a unit eigenvector $v$. So $\langle x, (M - I)x\rangle \ge 0$ for
every $x$, that is, $M - I$ is positive semidefinite, and $Mv = v$. Then $P = w(M - I)$ is positive
semidefinite, $\lambda_{\min}(wM) = w\lambda_{\min}(M) = w$, and $Px = 0$ exactly when $Mx = x$, as
$w > 0$: the null space $\ker P$ is the eigenspace $E_1$ of $M$ for $1$. For the range: if $Mx = cx$ with
$c > 1$ then $x = P\bigl((w(c-1))^{-1}x\bigr)$, so every eigenspace $E_c$, $c > 1$, lies in the range,
and so does their span $K$. The range is orthogonal to $E_1$: $P$ is symmetric, so
$\langle z, Px\rangle = \langle Pz, x\rangle = 0$ for $z \in E_1$. The eigenspaces of
the symmetric matrix $M$ span $\RR^d$, and $E_c = 0$ for $c < 1$, since
$0 \le \langle x, (M - I)x\rangle = (c - 1)|x|^2$ on $E_c$. So every $y$ is $a + k$ with $a \in E_1$ and
$k \in K$. If $y$ is in the range, then $a = y - k$ is in the range as well as in $E_1$, so
$\langle a, a\rangle = 0$, and $y = k \in K$: the range is $K$. By rank--nullity the rank of $P$ is $d$
minus the dimension of its null space, which contains $v \ne 0$, so $\rank P \le d - 1$. It is
$0$ exactly when $P = 0$, that is $M = I$. For $d = 2$ and $M \ne I$ it is therefore $1$.

\emph{(2).} By \ref{prop:affine-orbits}(1), $\widehat\mu_\Sigma(\omega) = e^{-g(\langle\omega,\Sigma\omega\rangle)}$.
If $\mu_{\Sigma_w} = \mu_{wI} * \rho$, taking transforms gives
$e^{-g(\langle\omega,(wI + P)\omega\rangle)} = e^{-g(\langle\omega,wI\omega\rangle)}\,\widehat\rho(\omega)$, which is
\ref{eq:one-step} with $\Sigma = wI \in \PD(d)$ and $P$ positive semidefinite with $0 < \rank P < d$ by
(1), and $g$ is continuous with $g(0) = 0$ and $g \ge 0$. \ref{thm:rank-deficient-step} gives $g(u) = bu$ with
$b = g(1) \ge 0$ and $\rho$ the centred Gaussian law with covariance $2bP$, which is carried by the range
of $P$, the law of $(2bP)^{1/2}Z$ with $Z$ standard normal. Then $\widehat\mu_I(\omega) = e^{-b|\omega|^2}$,
the transform of the centred Gaussian law with covariance $2bI$. Conversely, if $g(u) = bu$ with
$b \ge 0$, the converse half of \ref{thm:rank-deficient-step} gives \ref{eq:one-step} for $\Sigma = wI$ and
this $P$ with that Gaussian $\rho$, and the transforms of $\mu_{\Sigma_w}$ and $\mu_{wI} * \rho$ agree.

\emph{(3).} Let $(\mu_{\Sigma,\Sigma'})$ be a $(\GL(d), W)$-cascade family of kernels, with kernels from scale zero $\mu_\Sigma$. $W$ contains the rays, so the ray order is contained in it, and
\ref{lem:gw-kernels-from-origin}(1) restricts the family to $(\GL(d), \text{rays})$ with the same kernels
from scale zero. By \ref{prop:gw-ray-families}(3), $\mu_\Sigma = (\Sigma^{1/2})_*\mu_I$ with $\mu_I$
rotation invariant, self-decomposable and not $\delta_0$, and $\widehat\mu_\Sigma(\omega) =
F(\langle\omega,\Sigma\omega\rangle)$ with $F(u) := \widehat\mu_I(\sqrt u\,e_1)$ continuous and $F(0) = 1$. By
\ref{prop:gw-ray-families}(1), $F$ has no zero. It is real, as $\mu_I$ is invariant under $-I$, and
$|F| \le 1$, so $F > 0$, and $g := -\log F$ is continuous, nonnegative, with $g(0) = 0$ and
$e^{-g(|\omega|^2)} = \widehat\mu_I(\omega)$: the hypotheses of \ref{prop:affine-orbits} hold, and the
$\mu_\Sigma$, $\Sigma \in \PD(d)$, are the orbit of $\mu_I$. (K1) of \ref{def:gw-cascade-family} at $0 \le_W wI \le_W wM$ gives
$\mu_{wM} = \mu_{wI} * \mu_{wI,wM}$, so (2) gives $g(u) = bu$, and $b > 0$ because $\mu_I \ne \delta_0$.
Then $\widehat\mu_\Sigma(\omega) = e^{-b\langle\omega,\Sigma\omega\rangle}$ for every $\Sigma \in \PSD(d)$, by
the form above. These transforms have no zero, and (K1) at $0 \le_W \Sigma \le_W \Sigma'$ gives
$\widehat\mu_{\Sigma'} = \widehat\mu_{\Sigma,\Sigma'}\,\widehat\mu_\Sigma$, so
$\widehat\mu_{\Sigma,\Sigma'}(\omega) = e^{-b\langle\omega,(\Sigma'-\Sigma)\omega\rangle}$.
\end{proof}

\medskip

\textbf{The families across the rungs.} The draft's Remark~2.17 reads \ref{lem:loewner-scale-shape}
for a vision application, and it stays prose. Read in the nodes' terms, beyond the step of
\ref{prop:shape-adapted-apertures}, it is a ladder of four rungs. The Gaussian cascades in every
direction of $(t,S)$ (\ref{thm:affine-rigidity}). The isotropic stable orbits cascade across shapes
in any orientation, with an anisotropy gained per step of at most $(t_2/t_1)^{d\alpha/(d+\alpha)}$
(\ref{prop:stable-step-cone}), and so, by ordered \Levy{} densities on the same cone, do the orbits
of profiles that are stable at the origin and tempered beyond, which have exponential moments
(\ref{prop:tempered-step-cone}). A frame family cascades in scale and in relative scale at a fixed
orientation (\ref{prop:frame-families}), and the families of one linear flow drift in shape along a
fixed curve, locked to scale (\ref{prop:one-parameter-families}). The \Matern{} orbits, and
every orbit of logarithmic exponent, among them every orbit whose symbol is rational and does not
vanish, the kind elliptic solves in the shape's own operator implement exactly, cascade in scale only
(\ref{prop:matern-steps}, \ref{prop:log-orbit-rays}, \ref{prop:rational-symbols}), so for them a different shape is a re-measurement of the original
image rather than a further smoothing of the current one; between these ends, the steps an orbit
survives by ordered \Levy{} densities are read off its profile's dilation ratio
(\ref{prop:orbit-step-cone}). The draft's Remark~2.12 says the re-measurement of every non-Gaussian
aperture; that sentence is narrowed to the orbits of logarithmic exponent, where it is proved. For
families covariant under similarities only, an elliptical part at \Levy{} index $0$ is shape-blind
(\ref{prop:similarity-orbit-family}(3)--(4)), and a general part is shape-blind at the origin and need
not be beyond it (\ref{prop:similarity-general-index-zero}, \ref{prop:similarity-line-family}); at a
positive index the shape of the part is bounded when its rank-one limit is comparable to the isotropic
member (\ref{prop:similarity-general-bounded-shape}) and may be unbounded otherwise
(\ref{prop:similarity-degenerate-family}).
